\documentclass[10pt]{article}
\usepackage[margin=0.75in]{geometry}
\usepackage{amsmath,amssymb,amsthm,microtype,hyperref}
\newtheorem{theorem}{Theorem}
\begin{document}
\begin{center}
{\Large Arithmetic Applicability of Classical Odd-Run Merges}\par\medskip
Collatz proof project\quad---\quad September 5, 2026
\end{center}
This note formalizes the classical consecutive-seed endpoint identities
and applies them to the current proof program. The merging stem is due
to Garner; the explicit endpoint dichotomy appears in LaDue's Lemma 3.1
and Theorem 4.1\ [1]. No mathematical novelty or global coverage is claimed.
Throughout, $T(n)=n/2$ for even $n$ and $T(n)=(3n+1)/2$ for odd $n$.

\begin{theorem}[Odd-run endpoint dichotomy]
Let $a\geq1$, let $m$ be a positive odd integer, and put $n=2^am-1$.
The first $a$ steps from $n$ are odd, and
\[
T^a(n)+1=3^am,\qquad T^a(n)=3T^a(n-1)+2.
\]
At time $a+2$, exactly one of the following holds:
\[
\begin{array}{c|l}
3^am\pmod4&\text{endpoint relation}\\\hline
1&T^{a+2}(n)=T^{a+2}(n-1),\\
3&T^{a+2}(n)=9T^{a+2}(n-1)+2.
\end{array}
\]
In particular, equality at that time holds if and only if $3^am\equiv1\pmod4$.
\end{theorem}
\begin{proof}
For an odd input, $T(n)+1=3(n+1)/2$. Repeating this identity proves the
odd-run formula. The seed $n-1$ first takes an even step, followed by
$a-1$ odd steps, giving $T^a(n-1)+1=3^{a-1}m$. This proves the
ratio-three identity. Both resulting endpoints are even. If $3^am$ is
one modulo four, the path from $n$ takes two even steps and the path
from $n-1$ takes an even and an odd step; the endpoints agree.
If $3^am$ is three modulo four, these two patterns are reversed,
and direct substitution gives the ratio-nine identity. The latter
endpoints cannot agree for natural seeds.

In Lean, \texttt{realizes\_odd\_run} proves actual parity realization,
not just a formal affine formula. The arithmetic merge lemma
applies the existing stem theorem. The remaining branch and the equivalence
are proved separately by exact natural arithmetic.
\end{proof}

\begin{theorem}[Necessary exit for a least unbounded seed]
Suppose $n>1$ is the least positive seed with an unbounded orbit, and
$n+1=2^am$ with $a\geq1$ and $m$ odd. Then $3^am\equiv3\pmod4$.
Thus its initial odd run must be followed by a single even step and
then an odd step, rather than two even steps.
\end{theorem}
\begin{proof}
In the double-even case, the first theorem merges $n$ with $n-1$.
Minimality makes the positive seed $n-1$ bounded, so their common tail
is bounded, contradicting unboundedness of $n$.
No noncontraction property is needed for the replacement. This distinction
matters: $n-1$ is even and its first multiplicative coefficient is $1/2$.
The argument concerns the least unbounded seed, not the least seed whose
coefficients never contract.
\end{proof}

\paragraph{Breaking a finite-table limitation.}
The previous fixed-table obstruction uses seeds with
$n\equiv-1\pmod{2^{K+18}}$ and $n\equiv0\pmod{3^{11}}$.
For any $a\geq K+18$, choose the odd part $m$ to satisfy
\[
m\equiv2^{-a}\pmod{3^{11}},\qquad
m\equiv3^{-a}\pmod4.
\]
CRT gives infinitely many such $m$. The resulting seeds remain in the
table's obstruction through horizon $K$, yet the classical variable-length
stem merges them with $n-1$ after $a+2$ steps. Choosing instead
$m\equiv3\cdot3^{-a}\pmod4$ gives the single-even branch within the same
obstruction. The generic endpoint theorems are kernel checked; this CRT
application is a written deduction with exact Python instance checks.

\paragraph{What remains open.}
For a least unbounded $n$, the remaining branch gives $9z+2$ on its orbit
and $z$ on the bounded orbit of $n-1$. Nothing here proves that boundedness
of $z$ implies boundedness of $9z+2$. That missing implication, or a
different argument covering the branch, is substantive. Neither
nondivergence nor exclusion of nontrivial cycles follows from these results.

\paragraph{Verification and source.}
\begingroup\raggedright
Seven declarations in \texttt{Collatz.OddRunMerges} enter the axiom audit.
\texttt{analysis/odd\_run\_merges.py} constructs both branches, with tests
of actual orbits and table applicability. These are instance checks,
not a finite census claimed to prove a universal statement.\par\endgroup

\noindent[1] M. D. LaDue, \emph{Clusters of Integers with Equal Total Stopping
Times in the $3x+1$ Problem} (2017), Lemma 3.1 and Theorem 4.1.
\href{https://arxiv.org/abs/1709.02979}{arXiv:1709.02979}.
\end{document}
